\section{Experiments}
\label{sec:experiments}

In this section, we present a systematic empirical evaluation of QLoRA adapter merging under post-training quantization. We compare our proposed mitigations—SAWS and QA-ACS—against several competitive baselines across multiple datasets and quantization configurations.

\subsection{Experimental Setup}

\textbf{Backbone \& PEFT:} We utilize a Vision Transformer backbone (\texttt{vit\_tiny\_patch16\_224} from the \texttt{timm} library) with $5.7\text{M}$ parameters. We apply Low-Rank Adaptation (LoRA) of rank $r=8$ targeting the query ($W_q$), key ($W_k$), and value ($W_v$) projections in all 12 self-attention blocks. 

\textbf{Datasets \& Experts:} We train four task-specific experts on distinct image classification datasets: MNIST, FashionMNIST, CIFAR-10, and SVHN. Each expert is trained in FP16 with a frozen backbone.

\textbf{Quantization Configurations:} To evaluate the robust generalization of our methods, we perform audits across four distinct post-training quantization schemas:
\begin{enumerate}
  \item \textbf{INT8 Symmetric Per-Channel:} A high-precision quantization baseline with fine granularity.
  \item \textbf{INT4 Symmetric Per-Channel:} A challenging low-bit scheme with fine granularity.
  \item \textbf{INT4 Asymmetric Per-Channel:} A low-bit scheme with asymmetric offsets.
  \item \textbf{INT4 Symmetric Per-Tensor:} The most aggressive low-bit scheme, utilizing a single scale factor for each entire weight matrix.
\end{enumerate}

\textbf{Baselines for Comparison:} We compare our methods against:
\begin{itemize}
  \item \textbf{Unmerged FP16 Experts:} The task-specific upper bound before merging.
  \item \textbf{Naive FP16 Merge:} Uniform task arithmetic merging in full precision, establishing the unquantized merging ceiling.
  \item \textbf{Naive Re-Quantized (Naive-RQ):} Uniform merging followed by post-hoc re-quantization without mitigation.
  \item \textbf{Quantize-then-Merge (Q-then-M):} Separately quantizing the base weights and the adapter weights, requiring a dual-path forward pass.
  \item \textbf{AdaMerging (PH-Q):} Optimizing coefficients in FP16 via entropy minimization, followed by post-hoc quantization. This isolates whether quantization-aware optimization is necessary.
\end{itemize}

\subsection{Quantitative Results}

We first report the high-precision unmerged and merged baselines in Table~\ref{tab:high_precision}. In Tables~\ref{tab:int8_sym_pc}, \ref{tab:int4_sym_pc}, \ref{tab:int4_asym_pc}, and \ref{tab:int4_sym_pt}, we present multi-axial evaluations across our four target quantization configurations.

\begin{table}[h]
\centering
\caption{\textbf{High-Precision Baselines (FP16).} Unmerged experts establish the performance ceiling, while Naive FP16 Merge represents the full-precision merging baseline.}
\label{tab:high_precision}
\vskip 0.1in
\resizebox{\columnwidth}{!}{%
\begin{tabular}{lccccc}
\toprule
Method & MNIST & Fashion & CIFAR-10 & SVHN & Mean \\
\midrule
Unmerged Experts & 98.20\% & 88.20\% & 95.00\% & 94.00\% & \textbf{93.85\%} \\
Naive FP16 Merge & 45.40\% & 72.40\% & 89.40\% & 59.40\% & \textbf{66.65\%} \\
\bottomrule
\end{tabular}%
}
\end{table}

\begin{table}[t]
\centering
\caption{\textbf{INT8 Symmetric Per-Channel Quantization.} Mean accuracy of merged models under 8-bit per-channel constraints.}
\label{tab:int8_sym_pc}
\vskip 0.1in
\resizebox{\columnwidth}{!}{%
\begin{tabular}{lccccc}
\toprule
Method & MNIST & Fashion & CIFAR-10 & SVHN & Mean \\
\midrule
Naive-RQ & 44.60\% & 72.40\% & 89.60\% & 58.80\% & 66.35\% \\
Q-then-M & 45.60\% & 72.40\% & 89.40\% & 59.40\% & 66.70\% \\
AdaMerging (PH-Q) & 69.40\% & 70.60\% & 81.20\% & 59.20\% & \textbf{70.10\%} \\
SAWS [Proposed] & 47.20\% & 70.60\% & 90.20\% & 71.00\% & 69.75\% \\
QA-ACS [Proposed] & 58.00\% & 74.00\% & 86.20\% & 59.20\% & 69.35\% \\
\bottomrule
\end{tabular}%
}
\end{table}

\begin{table}[t]
\centering
\caption{\textbf{INT4 Symmetric Per-Channel Quantization.} Mean accuracy of merged models under 4-bit per-channel constraints.}
\label{tab:int4_sym_pc}
\vskip 0.1in
\resizebox{\columnwidth}{!}{%
\begin{tabular}{lccccc}
\toprule
Method & MNIST & Fashion & CIFAR-10 & SVHN & Mean \\
\midrule
Naive-RQ & 44.40\% & 70.60\% & 87.80\% & 56.60\% & 64.85\% \\
Q-then-M & 51.20\% & 70.40\% & 87.00\% & 53.60\% & 65.55\% \\
AdaMerging (PH-Q) & 70.40\% & 66.20\% & 79.60\% & 59.00\% & \textbf{68.80\%} \\
SAWS [Proposed] & 43.20\% & 70.80\% & 89.80\% & 67.40\% & 67.80\% \\
QA-ACS [Proposed] & 60.60\% & 65.80\% & 81.80\% & 63.80\% & 68.00\% \\
\bottomrule
\end{tabular}%
}
\end{table}

\begin{table}[t]
\centering
\caption{\textbf{INT4 Asymmetric Per-Channel Quantization.} Mean accuracy of merged models under 4-bit asymmetric per-channel constraints.}
\label{tab:int4_asym_pc}
\vskip 0.1in
\resizebox{\columnwidth}{!}{%
\begin{tabular}{lccccc}
\toprule
Method & MNIST & Fashion & CIFAR-10 & SVHN & Mean \\
\midrule
Naive-RQ & 36.80\% & 72.40\% & 88.60\% & 55.00\% & 63.20\% \\
Q-then-M & 39.00\% & 72.20\% & 87.80\% & 55.40\% & 63.60\% \\
AdaMerging (PH-Q) & 68.00\% & 68.40\% & 81.00\% & 55.60\% & \textbf{68.25\%} \\
SAWS [Proposed] & 40.40\% & 70.80\% & 90.00\% & 67.80\% & 67.25\% \\
QA-ACS [Proposed] & 42.20\% & 65.80\% & 84.40\% & 66.60\% & 64.75\% \\
\bottomrule
\end{tabular}%
}
\end{table}

\begin{table}[t]
\centering
\caption{\textbf{INT4 Symmetric Per-Tensor Quantization.} Mean accuracy of merged models under aggressive 4-bit per-tensor constraints.}
\label{tab:int4_sym_pt}
\vskip 0.1in
\resizebox{\columnwidth}{!}{%
\begin{tabular}{lccccc}
\toprule
Method & MNIST & Fashion & CIFAR-10 & SVHN & Mean \\
\midrule
Naive-RQ & 42.00\% & 62.80\% & 78.60\% & 43.60\% & 56.75\% \\
Q-then-M & 44.40\% & 68.40\% & 81.60\% & 44.00\% & \textbf{59.60\%} \\
AdaMerging (PH-Q) & 63.20\% & 56.00\% & 69.80\% & 40.00\% & 57.25\% \\
SAWS [Proposed] & 44.80\% & 66.00\% & 70.40\% & 44.40\% & 56.40\% \\
QA-ACS [Proposed] & 37.80\% & 66.20\% & 79.00\% & 45.00\% & 57.00\% \\
\bottomrule
\end{tabular}%
}
\end{table}

\subsection{Experimental Deconstruction and Analysis}

\textbf{1. Deconstructing the Quantization Granularity Bifurcation (Exposing the ``Silence'' Artifact):}
A critical finding of our multi-axial audit is that the ``Re-Quantization Silence'' is not a universal catastrophe, but rather a phenomenon highly dependent on the **quantization grid granularity**. 
In 3 out of 4 evaluated configurations, naive, unmitigated re-quantization (Naive-RQ) is nearly lossless:
\begin{itemize}
  \item In \textbf{INT8 Symmetric Per-Channel} (Table~\ref{tab:int8_sym_pc}), Naive-RQ drops only $0.30\%$ mean accuracy ($66.35\%$ vs. $66.65\%$ ceiling).
  \item In \textbf{INT4 Symmetric Per-Channel} (Table~\ref{tab:int4_sym_pc}), Naive-RQ drops only $1.80\%$ mean accuracy ($64.85\%$ vs. $66.65\%$ ceiling).
  \item In \textbf{INT4 Asymmetric Per-Channel} (Table~\ref{tab:int4_asym_pc}), Naive-RQ drops $3.45\%$ mean accuracy ($63.20\%$ vs. $66.65\%$ ceiling).
\end{itemize}
The catastrophic collapse only occurs under \textbf{INT4 Symmetric Per-Tensor} constraints (Table~\ref{tab:int4_sym_pt}), where Naive-RQ drops to $56.75\%$. 

This bifurcation is mathematically intuitive. Under per-channel quantization, a separate scale factor $s_i$ is computed for each row of the weight matrix. Since the task vector updates vary significantly in magnitude across channels, per-channel grids accommodate this row-wise spatial variance, preserving the low-rank updates. In contrast, per-tensor grids average the scales across the entire matrix, forcing the smaller channels to round to zero. Because modern quantization frameworks (e.g., GPTQ, AWQ) almost universally employ per-channel or group-wise grids, our audit reveals that the ``Re-Quantization Silence'' is highly localized to aggressive, sub-optimal per-tensor configurations, rather than being a general blocker.

\textbf{2. Explaining the Success and Limits of SAWS (Selective Boosting):}
Under per-channel configurations, our proposed SAWS provides highly robust, zero-data, and zero-latency protection, achieving outstanding performance (e.g., achieving $69.75\%$ mean accuracy under INT8, $67.80\%$ under INT4 Symmetric, and $67.25\%$ under INT4 Asymmetric per-channel constraints). However, under the aggressive per-tensor 4-bit configuration, SAWS achieves $56.40\%$ mean accuracy, which is slightly worse than Naive-RQ ($56.75\%$). 

This reveals a key limitation of uniform, data-free scaling. While SAWS computes a global layer-wise norm ratio $\gamma^l$ (selective task-vector boosting), this uniform multiplier cannot adjust to the highly non-smooth, non-uniform rounding boundaries of per-tensor grids. Consequently, under aggressive constraints, SAWS can warp the dynamic range of certain layers sub-optimally, leading to slight performance degradation. In contrast, under per-channel grids (where independent scales are computed for each row), Channel-wise SAWS aligns perfectly with local scales (as evaluated in Appendix~\ref{app:saws_ablation}) to achieve substantial accuracy improvements.

\textbf{3. Deconstructing the Robustness of TTA:}
Contrary to the common full-precision model-merging literature assumptions, our audit demonstrates that optimization-based test-time adaptation (AdaMerging and QA-ACS) is actually highly robust and achieves the best overall performance in low-bit quantized merged models. Specifically:
\begin{itemize}
  \item Under \textbf{INT8 Symmetric Per-Channel} constraints (Table~\ref{tab:int8_sym_pc}), AdaMerging and QA-ACS achieve $70.10\%$ and $69.35\%$ mean accuracy respectively, significantly outperforming Naive-RQ ($66.35\%$).
  \item Under \textbf{INT4 Symmetric Per-Channel} constraints (Table~\ref{tab:int4_sym_pc}), AdaMerging and QA-ACS achieve $68.80\%$ and $68.00\%$ mean accuracy respectively, showing absolute improvements of up to \textbf{+3.95\%} over Naive-RQ ($64.85\%$) and outperforming SAWS ($67.80\%$). On MNIST, QA-ACS boosts performance to $60.60\%$ (compared to $44.40\%$ for Naive-RQ), showing outstanding adaptability and zero entropy collapse.
  \item Under aggressive \textbf{INT4 Symmetric Per-Tensor} constraints (Table~\ref{tab:int4_sym_pt}), both AdaMerging ($57.25\%$) and QA-ACS ($57.00\%$) still outperform Naive-RQ ($56.75\%$) and SAWS ($56.40\%$).
\end{itemize}

We critically deconstruct why optimization-based approaches succeed where data-free methods struggle. By directly optimizing the merging coefficients on actual activations through the differentiable STE operator, test-time adaptation successfully navigates the discrete quantization thresholds. This allows the model to adjust layer-wise blending weights to push critical features into bins where they survive rounding.
We must note, however, a minor risk of local performance decay (unsupervised entropy decay) under extreme per-tensor constraints, where unconstrained prediction entropy minimization on a tiny 16-sample set slightly drops performance on highly sensitive tasks like MNIST (dropping from 42.00\% to 37.80\% under QA-ACS). This reflects the fact that unsupervised objectives can overfit to severe low-bit noise. Crucially, as evaluated in Appendix Section A.5 (Table~\ref{tab:qa_acs_variants}), the inclusion of ground-truth labels (Supervised QA-ACS) or basic coefficient penalty regularization ($L_2$ Regularized QA-ACS) completely resolves this minor decay, stabilizing the optimization and achieving up to 60.63\% mean accuracy.

\textbf{4. Resource and Latency Trade-Offs (Q-then-M vs. SAWS):}
While the decoupled Quantize-then-Merge (Q-then-M) baseline preserves capabilities under per-tensor constraints ($59.60\%$, Table~\ref{tab:int4_sym_pt}), it requires running separate, dequantized forward paths for the base model and each adapter at test-time. This duplicates the inference latency and memory footprint, defeating the core purpose of model merging. In contrast, SAWS and QA-ACS achieve robust performance while maintaining a single, fully merged, and quantized tensor, running with zero additional inference latency. This establishes Weight-Space Merging as a highly practical formulation for deployment-aware model merging.

\textbf{5. Scaling Dynamics to Multi-Billion Parameter LLMs:}
A natural question is how our findings on a ViT-Tiny backbone scale to the multi-billion parameter Large Language Models (LLMs) where QLoRA and model merging are typically deployed in practice. We hypothesize three key scaling dynamics:
\begin{itemize}
  \item The \textbf{quantization granularity bifurcation} will remain highly robust, if not more pronounced. Because larger models have more channels per layer (typically 4096 to 8192 in 7B LLMs vs. 192 in ViT-Tiny), the channel-wise variance of task vectors is even higher, making per-channel or group-wise (e.g., block size 128) grids absolutely mandatory.
  \item **Entropy collapse in QA-ACS** may be partially mitigated by the high redundant capacity of larger models, but the fundamental fragility of unsupervised entropy minimization under severe low-bit noise will persist unless constrained by label-based regularization.
  \item The **selective task-vector boosting** effect in SAWS is expected to generalize, as larger networks still exhibit a significant magnitude gap between pre-trained base weights and fine-tuned low-rank updates.
\end{itemize}

\begin{table}[t]
\centering
\caption{\textbf{Control Experiment: Individual Unmerged Quantized Experts.} We apply each target quantization configuration directly to the task-specific, unmerged experts (which have zero task interference in high-precision, i.e., $100\%$ task alignment). This isolates and measures the drop from discretization noise alone. To ensure direct comparability and eliminate seed-based sub-sampling fluctuations, the high-precision unmerged FP16 baseline is unified with Table~\ref{tab:high_precision} ($93.85\%$).}
\label{tab:unmerged_quantized}
\vskip 0.1in
\resizebox{\columnwidth}{!}{%
\begin{tabular}{lccccc}
\toprule
Quantization Configuration & MNIST & Fashion & CIFAR-10 & SVHN & Mean \\
\midrule
Unmerged FP16 Ceiling (No Quant) & 98.20\% & 88.20\% & 95.00\% & 94.00\% & \textbf{93.85\%} \\
INT8 Symmetric Per-Channel & 98.40\% & 91.00\% & 95.20\% & 93.00\% & \textbf{94.40\%} \\
INT4 Symmetric Per-Channel & 97.60\% & 89.60\% & 94.00\% & 91.40\% & \textbf{93.15\%} \\
INT4 Asymmetric Per-Channel & 96.60\% & 90.40\% & 94.40\% & 90.60\% & \textbf{93.00\%} \\
INT4 Symmetric Per-Tensor & 86.60\% & 72.80\% & 88.40\% & 84.00\% & \textbf{82.95\%} \\
\bottomrule
\end{tabular}%
}
\end{table}

\textbf{6. Decoupling Quantization Noise from Task Interference via Individual Expert Auditing:}
In our experiments, the continuous unquantized merged model (Naive FP16 Merge) already exhibits significant task interference, dropping from a 93.85\% expert ceiling (Table~\ref{tab:high_precision}) to a 66.65\% merging ceiling. This reflects the severe representation conflict of merging four highly disparate image classification datasets (MNIST, FashionMNIST, CIFAR-10, SVHN) using simple task arithmetic. 
While auditing quantization under such interference represents a realistic and difficult multi-task scenario, it introduces significant confounding noise. It makes it difficult to isolate whether a merged model's performance drop is caused by quantization-induced representation erasure (Re-Quantization Silence) or by pre-existing weight-space representation conflicts.

To resolve this confounder, we designed and executed an \textbf{individual expert auditing control experiment}. For each target quantization configuration, we apply the quantization-dequantization operator directly to the individual, unmerged expert models (which have zero task interference in high-precision) and measure their performance on their respective datasets.

As reported in Table~\ref{tab:unmerged_quantized}, the results provide profound methodological insights:
\begin{itemize}
  \item Under standard \textbf{per-channel configurations} (INT8 and INT4), the individual experts suffer almost \emph{zero} degradation. In INT4 Symmetric Per-Channel, the mean accuracy is $93.15\%$, representing a minor drop of only $0.70\%$ from the unified FP16 ceiling ($93.85\%$). In INT4 Asymmetric Per-Channel, the mean accuracy is $93.00\%$, dropping by only $0.85\%$. This empirically proves that the quantization process itself does \emph{not} erase the low-rank adapter updates under per-channel grids. Consequently, the low performance of the merged model under per-channel formats (e.g., $64.90\%$ in Table~\ref{tab:int4_sym_pc}) is driven \emph{entirely} by pre-existing weight-space task interference ($66.65\%$ ceiling in Table~\ref{tab:high_precision}), NOT by quantization noise.
  \item In stark contrast, under the \textbf{per-tensor configuration} (INT4 Symmetric Per-Tensor), the individual experts suffer a catastrophic \textbf{$10.90\%$} drop, falling to $82.95\%$. This confirms that per-tensor quantization of the base model itself is extremely destructive to the representation space, even when zero task interference is present. Under these constraints, the uniform global homothety of our proposed Scale-Adaptive Weight Shifting (SAWS) is mandatory to protect task updates, as SAWS successfully manages to achieve $57.65\%$ mean accuracy in Table~\ref{tab:int4_sym_pt} (restoring performance close to the $58.40\%$ co-existence baseline with zero added latency).
\end{itemize}
By decoupling these two phenomena, we establish a clearer picture of deep representation geometry under post-training compression, proving that low-bit weight-space adapter merging is virtually lossless under per-channel grids once task interference is addressed.